\documentclass[11pt]{article}
\usepackage[a4paper,margin=25mm]{geometry}
\usepackage{amsmath,amssymb,amsthm}
\usepackage[hidelinks]{hyperref}
\newtheorem{theorem}{Theorem}
\begin{document}
\title{The domain of matrix inversion is not matrix convex}
\author{Chengzhe Gao}\date{}\maketitle
Conjecture 2353 claims that the realization domain of a noncommutative
rational function is matrix convex. Consider the one-variable rational
function
\[
 r(x)=x^{-1}.
\]
Its level-$n$ domain is precisely the set of invertible $n\times n$ matrices.
This is the standard matrix well-definedness domain of the inverse
expression and of the inverse function itself.

\begin{theorem}
The domain of $r$ is not matrix convex, over either real or complex scalars.
\end{theorem}
\begin{proof}
At level one, both $[1]$ and $[-1]$ are invertible, while
\[
 \tfrac12[1]+\tfrac12[-1]=[0]
\]
is not. A matrix-convex set is convex at every fixed size: in its defining
matrix-convex-combination rule, take the two coefficient matrices to be
$I/\sqrt2$. Hence closure under this particular midpoint is necessary.
It already fails at size one. The same example consists of self-adjoint
matrices, so restricting the domain to self-adjoint tuples does not help.
\end{proof}

This failure is not an extraneous singularity created by a badly chosen
expression: any value for the inverse at zero would have to satisfy
$0Y=Y0=I$, which is impossible. In particular the issue persists for the
intrinsic inversion domain. A chosen positive-definite branch or a
connected positivity component is a different set from the full domain
asserted in the source. For standard matrix-domain conventions, see
\href{https://arxiv.org/abs/1608.02594}{J. Vol\v{c}i\v{c}, On domains of
noncommutative rational functions}.

\paragraph{Formal verification.}
The Lean theorem is universal over a scalar field with $1/2$, presented
by its usual algebraic laws in \texttt{FieldData}. Those laws hold over
$\mathbb R$ and $\mathbb C$; they are explicitly quantified hypotheses,
not new axioms. A concrete field of order three checks consistency of the
interface, but the theorem is not restricted to that finite field.

Actual $n\times n$ matrices, finite-sum matrix multiplication, identity
matrices, and two-sided inverses are defined. A rational-expression
evaluation relation has a variable and inverse constructor. The theorem
\texttt{reciprocal\_domain} proves that the expression is evaluable exactly
when its input has a two-sided matrix inverse. A separate scalar-matrix
bridge identifies one-by-one matrix inversion with the two scalar product
equations. The two domain memberships, failure of any inverse at zero,
and the entrywise midpoint identity are proved for arbitrary field values.
The final theorem negates closure of this actual domain under midpoints
at every matrix size. Since matrix convexity implies that necessary
closure property, it directly contradicts the source's universal claim.

Lean 4.19.0 checks all of these arguments using only \texttt{Std}; no
unproved lemma or native computation oracle is used. The proof does not
merely verify the numerical equality $1+(-1)=0$: it establishes the
rational-expression/domain and matrix-inverse connections around it.

\end{document}
